\chapter{Excessive Thirst (Polydipsia)}

\section*{Definition}
Polydipsia is the pathological sensation of excessive thirst leading to increased fluid intake, defined as consumption exceeding 3 liters per day in adults. It represents a symptom of underlying homeostatic derangement rather than a disease itself.

\section*{Pathophysiology}
Thirst is regulated by osmoreceptors in the hypothalamus that sense plasma osmolality and baroreceptors that detect volume depletion. Increased osmolality or decreased blood volume stimulates the thirst center, triggering antidiuretic hormone (ADH) release from the posterior pituitary. Polydipsia arises when osmotic or volume-regulating mechanisms are persistently activated or dysregulated.

\section*{Causes}
\begin{itemize}
  \item Diabetes mellitus (hyperglycemia-induced osmotic diuresis -- most common cause)
  \item Diabetes insipidus (central or nephrogenic -- inability to concentrate urine)
  \item Primary polydipsia (psychogenic or dipsogenic -- excessive drinking due to hypothalamic dysfunction)
  \item Hypercalcemia (elevated calcium impairs renal concentrating ability)
  \item Hypokalemia (potassium depletion impairs renal concentrating mechanism)
  \item Medications: Anticholinergics, diuretics, lithium (nephrogenic DI), phenytoin
  \item Renal disease (chronic kidney disease, post-obstructive diuresis)
  \item Sickle cell disease (renal medullary damage)
  \item Excessive sweating, diarrhea, vomiting (volume depletion)
\end{itemize}

\section*{When to See a Doctor}
See your doctor if:
\begin{itemize}
  \item Polydipsia persisting beyond 48 hours without obvious cause
  \item Associated polyuria (excessive urination), especially at night (nocturia)
  \item Unexplained weight loss with polydipsia (suggests diabetes mellitus)
  \item Headache, visual changes, or fatigue (may indicate diabetes insipidus or hypercalcemia)
  \item Concurrent nausea, vomiting, or abdominal pain (suggests hypercalcemia)
  \item Infants or elderly with polydipsia (higher risk of dehydration)
\end{itemize}

\section*{Related Symptoms}
\begin{itemize}
  \item Polyuria
  \item Nocturia
  \item Fatigue
  \item Weight loss
  \item Dry mouth
  \item Blurred vision
  \item Headache
  \item Nausea
\end{itemize}

\section*{Self-Care / Home Management}
\begin{itemize}
  \item Monitor fluid intake and urine output to quantify severity.
  \item Maintain hydration with water; avoid sugary drinks that may worsen hyperglycemia.
  \item Check capillary blood glucose if diabetic or at risk.
  \item Avoid excessive caffeine or alcohol, which have diuretic effects.
  \item Seek medical evaluation for persistent thirst, especially with polyuria or weight loss.
\end{itemize}

\section*{How Your Doctor Will Evaluate You}
\begin{itemize}
  \item Fasting blood glucose and HbA1c to screen for diabetes mellitus.
  \item Serum electrolytes, calcium, and renal function panel.
  \item Urinalysis with specific gravity and osmolality.
  \item Water deprivation test if diabetes insipidus suspected.
  \item MRI brain if central diabetes insipidus or hypothalamic lesion suspected.
\end{itemize}

\section*{Treatment Options}
\begin{itemize}
  \item Treat underlying cause (insulin for diabetes, desmopressin for central DI, hydration for hypercalcemia).
  \item Behavioral interventions for primary polydipsia (fluid restriction, cognitive behavioral therapy).
  \item Medication adjustment if drug-induced (e.g., lithium-induced nephrogenic DI).
  \item Electrolyte monitoring and repletion as needed.
\end{itemize}

\clearpage
